\section{Thin transverse strips at an arbitrary collision mark}
\label{sec:v48-thin-strips}
The endpoint envelopes of Section~\ref{sec:v47-endpoint-decisions}
are taken at fixed depth. They do not give a bound for a decision whose
depth increases with the collision count. We obtain that bound by
inserting a thin strip at its actual collision time. Its strong
multiplier norm is bounded independently of the strip width, while its
strong-to-weak norm tends to zero. The full occupation spectrum then
converts this smallness into an exact-coefficient local upper bound.
No long pullback of a strip is used as a multiplier.

The stable length exponent in \eqref{eq:v35-stable-norm} is
$\varsigma=1/8$. In this section put
\begin{equation}\label{eq:v48-strip-exponent}
 \alpha_0=\varsigma/2=1/16.
\end{equation}
The geometric number $q_0=47/53$ used below is unrelated to the
stable test exponent denoted by $q$ in the collision norms.

\subsection{Uniform multipliers and a small weak norm}
For $0<s<s_0$, let $b_{R,s}$ be the indicator of the union of the
coordinate strips of width $2s$ about the supporting lines of all
sides of $Y_R^*$. Angular strips are interpreted in fixed circle
charts. The horizontal strips lie in a fixed compact set away from
grazing. We may increase their widths by any one fixed factor without
changing the assertions. In particular
\[
 \one_{\{\operatorname{dist}(x,\partial Y_R^*)<s\}}\le b_{R,s}(x).
\]
The strips may overlap. Their number is fixed, not proportional to a
collision count. Define $b_{R,0}=0$ modulo the null supporting lines.

\begin{lemma}[A thin strip as a strong-to-weak multiplier]
\label{lem:v48-thin-multiplier}
On the local common collision spaces, uniformly in $R$ and $s$,
\begin{align}
 \|M_{b_{R,s}}h\|_{\mathcal B}&\le C\|h\|_{\mathcal B},
                                                        \notag\\
 |M_{b_{R,s}}h|_w&\le Cs^{\varsigma}\|h\|_{\mathcal B}.
                                      \label{eq:v48-strip-norms}
\end{align}
The first estimate does not assert smallness in the strong operator
norm. Both assertions also hold for any one of the strips and for
unions of a fixed bounded number of such strips.
\end{lemma}
\begin{proof}
Partition by all strip endpoints in $(\alpha,p)$ coordinates. These
are vertical or horizontal lines; the latter remain horizontal in
$(\alpha,\varphi)$. There is a common finite number of partition
cells in each homogeneity strip. A homogeneous stable graph has
slope bounded away from zero and infinity. It meets a cell in a
bounded number of intervals and spends length at most $C\epsilon$
within distance $\epsilon$ of its boundary. These constants do not
depend on the separation between two parallel partition lines.
Lemma~3.3(b) of the pinned version of \cite{DPZ}, on page 13,
therefore gives the first estimate, as in
Lemma~\ref{lem:v37-section-multiplier}.

For clarity, the uniformity as a strip collapses can also be seen in
the norm calculation. In the stable norm, splitting a test on $W$
into at most $N$ subintervals costs at most $N$ in the sum of
$(|U|/|W|)^{\varsigma}$. In the unstable comparison, matching the
strip endpoints removes intervals of length at most $C\epsilon$.
If the strip width is smaller than the matching distance, its whole
intersection may be assigned to that unmatched part and has the same
bound. Its stable-norm cost is $C\epsilon^{\varsigma}$, which is
admissible since $\gamma<\varsigma$. On the retained pieces the
indicator is constant. No inverse strip width enters either bound.
Membership in the strong completion and extension from smooth
densities follow from the same piecewise multiplier lemma.

To prove the second assertion, start with a smooth density $h$ and a
weak test $\psi$ on $W$. The intersection of $W$ with the strip union
is a union of at most $N$ intervals $U$, each of length at most $Cs$.
Restriction of a $C^p$ test to $U$ has uniformly bounded $C^q$ norm.
The length normalization of the strong stable norm gives
\[
 \left|\int_U h\psi\,dm_U\right|
       \le C\|h\|_s |U|^{\varsigma}.
\]
Summing these finitely many intervals and taking the weak supremum
proves the bound. Strong approximation and boundedness of the
multiplier extend it to $h\in\mathcal B$. At $s=0$ the same bound
gives the zero operator. The argument is uniform in the locations of
the lines, hence also in the moving radius.
\end{proof}

\subsection{Weak smallness on a moving spectral chart}
Write $Q=\mathscr T_{R,\theta}$ for the exact uncentered full-occupation
operator of \eqref{eq:v40-physical-operator}. Fix the finitely many
peak charts in Lemma~\ref{lem:v41-moving-maxima}. They all lie in
one fixed small roof neighborhood, independently of a later outer
roof band. On a smaller chart write
\[
 Q^a=\lambda^a\Pi+N^a,\qquad
 \|N^a\|\le C\rho^a,
 \qquad \|Q^a\|\le C \quad(a\ge0).
\]
Here and below $N^0$ in this formula means $I-\Pi$.
The last bound follows from $|\lambda|\le1$, the bounded projection,
and the complementary estimate. Use a common equivalent norm on
this fixed finite cover. The full Doeblin--Fortet estimate gives
\begin{equation}\label{eq:v48-chart-DF}
 \|Q^a h\|\le C\sigma^a\|h\|+C(1+a)|h|_w,
 \qquad |Q^a h|_w\le C(1+a)|h|_w,
 \qquad \sigma<1.
\end{equation}
Shrink the peak charts so that $|\lambda|\ge r>\sigma^{1/4}$.
This is possible since their central eigenvalues have modulus one.
The constants in this paragraph concern the peak neighborhood only.

\begin{lemma}[Projection and terminal-function interpolation]
\label{lem:v48-weak-interpolation}
Uniformly on these charts,
\begin{equation}\label{eq:v48-projection-interpolation}
 \|\Pi h\|+\sup_{a\ge0}|\ell(Q^a h)|
               \le C\|h\|^{1/2}|h|_w^{1/2}.
\end{equation}
Consequently, for every $a\ge0$,
\begin{align}
 \|\Pi M_{b_{R,s}}Q^a\nu\|&\le Cs^{\alpha_0},\notag\\
 |\ell(Q^aM_{b_{R,s}}\Pi\nu)|&\le Cs^{\alpha_0}.
                                      \label{eq:v48-small-amplitudes}
\end{align}
\end{lemma}
\begin{proof}
By homogeneity take $\|h\|=1$ and put $t=|h|_w$.
The case $t$ bounded below follows from power boundedness. For
$0<t<1/2$, choose an integer $L$ with $\sigma^L\le t<\sigma^{L-1}$.
Commutation of $Q$ and $\Pi$ gives
\[
 \|\Pi h\|=|\lambda|^{-L}\|\Pi Q^Lh\|
       \le C r^{-L}(\sigma^L+(1+L)t)
       \le C(1+|\log t|)t^{3/4}\le Ct^{1/2}.
\]
We used $-\log r/(-\log\sigma)<1/4$. For $a\le L$, continuity of
mass on the weak space and \eqref{eq:v48-chart-DF} give
$|\ell(Q^a h)|\le C(1+L)t$. For $a>L$, first apply the strong
bound in \eqref{eq:v48-chart-DF} at time $L$, then the uniform strong
power bound for $Q^{a-L}$. This gives the same upper bound, with a
changed constant. It is at most $Ct^{1/2}$.
If $t=0$, let $L\to\infty$ in these estimates: the projection is
zero because $\sigma<r$, and every fixed terminal pairing is zero by
the weak inequality. Restoring homogeneity proves
\eqref{eq:v48-projection-interpolation}.
Apply it to $M_{b_{R,s}}Q^a\nu$ and $M_{b_{R,s}}\Pi\nu$.
Their strong norms are bounded; their weak norms are at most
$Cs^{\varsigma}$ by Lemma~\ref{lem:v48-thin-multiplier}.
This proves \eqref{eq:v48-small-amplitudes}.
\end{proof}

\begin{theorem}[A local upper bound uniform in strip width and mark]
\label{thm:v48-marked-strip-local}
There is $C<\infty$ with the following property. For each fixed
$B\ge B_0$ there are $C_B<\infty$ and $\rho_B<1$ such that for every
$m\ge2$, $0\le j\le m$, $0<s<s_0$, radius $R$, discrete targets
$k\in\Z^2$, $l\in\Z$, real $t$, and interval $I$ of length $h>0$,
\begin{align}
 &m^2\nu\{K^{\rm c}_{m,R}=k,A_{m,R}=l,S_{m,R}-t\in I,
                         b_{R,s}(T_R^jx)=1\}\notag\\
 &\hspace{12mm}\le C s^{\alpha_0}(h+B^{-1})
                  +C_B(1+Bh)m^2\rho_B^{m/2}.
                                      \label{eq:v48-marked-strip-local}
\end{align}
The leading constant is independent of the outer band $B$. In
particular $s$ and $j$ may depend on $m$; the error does not depend
on them. This is an upper bound, not a Gaussian-amplitude theorem
for a shrinking strip.
\end{theorem}
\begin{proof}
The exact characteristic pairing with the mark is
\begin{equation}\label{eq:v48-marked-pairing}
 \ell(Q^{m-j}M_{b_{R,s}}Q^j\nu)
   =\int b_{R,s}(T_R^jx)
      e^{i\theta\cdot(K^{\rm c}_{m,R},S_{m,R},A_{m,R})}\,d\nu(x).
\end{equation}
This follows first on physical densities from
\eqref{eq:v40-physical-operator}. The sharp strip is the bounded
piecewise multiplier of Lemma~\ref{lem:v48-thin-multiplier}; its
physical multiplication agrees with that realization. One can check
the last assertion by smooth monotone averages of parallel-strip
indicators: their operator norms are uniformly bounded and their
strong-to-weak differences tend to zero, while their physical
integrals converge by dominated convergence. Thus no orbit-length
regularity of $b_{R,s}\circ T_R^j$ is needed.

On a peak chart, if $m-j\ge m/2$, expand the left long power in
\eqref{eq:v48-marked-pairing} and use the first inequality of
\eqref{eq:v48-small-amplitudes}. If $j\ge m/2$, expand the right
long power and use its second inequality. The unexpanded short power
is uniformly bounded in the strong norm. In either case
\[
 |\ell(Q^{m-j}M_{b_{R,s}}Q^j\nu)|
       \le Cs^{\alpha_0}|\lambda|^{m/2}+C\rho^{m/2}.
\]
The peak bound \eqref{eq:v41-peak-bound} integrates the first term
in four real frequency coordinates to $Cs^{\alpha_0}m^{-2}$.
On the nonresonant part of any fixed outer band, expand a long power
using the uniform spectral gap. The other power and the strip
multiplier remain bounded, giving $C_B\rho_B^{m/2}$.
These assertions hold on a finite parameter and frequency cover.

Choose the positive upper envelope of $\one_I$ with Fourier support
in $[-B,B]$, as in Section~\ref{sec:collision-mixing-local-limit}.
It is nonnegative, has mass $h+2\pi/B$, and its Fourier transform
has supremum at most that mass. Positivity gives the upper comparison
before Fourier inversion. Torus orthogonality, with the same signs
as the original fixed-coefficient inverse, and the preceding
integrated estimate prove \eqref{eq:v48-marked-strip-local}.
The nonpeak integral is at most a constant times
$B(h+B^{-1})\rho_B^{m/2}$. All actual resonances have roof coordinate
zero, so the peak cover and its leading constant are chosen once,
independently of $B$. No arithmetic residue has been discarded.
\end{proof}

\begin{corollary}[Summation over a graded family of marks]
\label{cor:v48-graded-marks}
Let $\mathcal J\subset\{0,\ldots,m\}$ and choose widths $s_j<s_0$.
The left side of \eqref{eq:v48-marked-strip-local}, with the strip
condition replaced by its union over $j\in\mathcal J$, is at most
\[
 C(h+B^{-1})\sum_{j\in\mathcal J}s_j^{\alpha_0}
       +C_B(1+Bh)|\mathcal J|m^2\rho_B^{m/2}.
\]
This finite-count inequality may be summed before taking a limit.
\end{corollary}
\begin{proof}
Apply the union bound to the positive, unchanged coefficient and use
Theorem~\ref{thm:v48-marked-strip-local} at every mark. The uniform
error, rather than an interchange of separate limsups, justifies the
assertion when the number or locations of the marks depend on $m$.
\end{proof}
